
The execution model defined by the \openshmem specification is extended by the
introduction of a \ac{GPU} symmetric heap allocated from GPU-attached memory.
This \ac{GPU} symmetric heap is distinct from, and in addition to, the default
symmetric heap defined by the \openshmem specification.

\subsection{Processing Elements and GPU Association}\label{subsec:pe_gpu_association}

The concept of a processing element (\ac{PE}) is defined by the \openshmem
specification. In this version of the specification, every \ac{PE} in an
\openshmem program that uses the GPU support defined herein shall be associated
with exactly one \ac{GPU}. The association between a \ac{PE} and its \ac{GPU}
shall be established during initialization of the \openshmem library and shall
remain fixed for the duration of the program.

The mechanism by which a \ac{GPU} is selected for association with a \ac{PE} is
outside the scope of this specification. The application or runtime environment
shall select the \ac{GPU} prior to initialization of the \openshmem library. The
\openshmem implementation shall use the \ac{GPU} that has been made current or
otherwise designated by the execution environment at the time of library
initialization.

\subsection{GPU Symmetric Heap}\label{subsec:gpu_symmetric_heap}

The \ac{GPU} symmetric heap is distributed across \acp{PE}, with each \ac{PE}'s
allocation residing in the memory of its associated \ac{GPU}. Routines are
provided for allocating and managing the \ac{GPU} symmetric heap memory from the host,
and for performing one-sided read (get), write (put), atomic, signaling, and synchronization
operations on symmetric objects within it from both the host and \ac{GPU}
execution contexts. These routines follow the \ac{PGAS} programming model
established by the \openshmem specification.

\subsection{Memory Domains}\label{subsec:memory_domains}

Because symmetric data may reside in either host-attached or GPU-attached memory,
this specification distinguishes two \emph{memory domains}. A memory domain is
the address space with respect to which the ordering, completion, and atomicity
of \openshmem operations are defined. Each \ac{PE} that uses the GPU support
defined herein exposes two distinct memory domains:

\begin{description}
\item[Default memory domain] The default symmetric heap, together with the
symmetric static and global data objects defined by the \openshmem
specification, residing in host (CPU-attached) memory.
\item[\ac{GPU} memory domain] The \ac{GPU} symmetric heap, residing in the
device (GPU-attached) memory of the \ac{GPU} associated with the \ac{PE}
(\cref{subsec:pe_gpu_association}).
\end{description}

Private (non-symmetric) host or device memory belongs to neither memory domain
and cannot be the remote target of an \openshmem operation.

The two memory domains are independent. Operations issued to one memory domain
are neither ordered nor completed with respect to operations issued to the other
memory domain unless a completion routine such as \FUNC{shmem\_quiet} is called,
or unless a routine explicitly provides cross-domain ordering, such as the
put-with-signal routines. The corresponding ordering and completion rules are
defined in \cref{subsec:gpu_aware_ordering} for GPU-aware operations and in
\cref{subsec:gc_ordering} for GPU-centric operations.

\textbf{Interaction with teams and atomicity domains:} Memory ordering and
completion are tracked per memory domain and per communication context, while
atomicity is defined per atomicity domain (\cref{subsec:atomicity_domains}). The
predefined device teams \CONST{SHMEM\_DEVICE\_TEAM\_WORLD} and
\CONST{SHMEM\_DEVICE\_TEAM\_SHARED} form atomicity domains that are independent
of those rooted at the host teams \CONST{SHMEM\_TEAM\_WORLD} and
\CONST{SHMEM\_TEAM\_SHARED}. The team \CONST{SHMEM\_DEVICE\_TEAM\_SHARED}
identifies the \acp{PE} that share a \ac{GPU} memory domain and can therefore
directly load/store one another's \ac{GPU} symmetric heap objects
(\cref{subsec:shmem_device_team_shared}). A communication context created with
the \CONST{SHMEM\_DEVICE\_CTX} option operates on the \ac{GPU} memory domain;
a context created without that option operates on the default memory domain.

\textbf{Memory domains targeted by each operation type:} The following table
summarizes which memory domains may be named by the operands of each class of
operation.

\begin{longtable}{|p{4.6cm}|p{9.9cm}|}
\hline
\textbf{Operation type} & \textbf{Memory domains that may be targeted} \\
\hline
\endfirsthead
\hline
\textbf{Operation type} & \textbf{Memory domains that may be targeted} \\
\hline
\endhead
CPU / host-initiated (\openshmem 1.6) & Default memory domain only. \\
\hline
GPU-aware (\cref{sec:gpu_aware}) & Default memory domain and \ac{GPU} memory
domain; a single host-initiated routine may name operands in either domain. \\
\hline
GPU-centric (\cref{sec:gpu_centric}) & \ac{GPU} memory domain only; source,
destination, and synchronization operands shall reside in the \ac{GPU} symmetric
heap. \\
\hline
\end{longtable}

